\appendix
\renewcommand{\thechapter}{\Alph{chapter}}
\titleformat{\chapter}[display]
  {\normalfont\centering\bfseries}
  {\Large APPENDIX \thechapter}{10pt}{\Large}
\setcounter{chapter}{0}

\chapter{Input Feature Schema}
\label{app:features}

Table~\ref{tab:app_schema} gives the complete input specification enforced by
the API and the web form: the accepted range, the reference value used by the
perturbation explanation, and the normal range displayed to the user.

\begin{table}[H]
\centering
\caption{Input feature schema with validation ranges and explanation reference
values.}
\label{tab:app_schema}
\small
\begin{tabular}{@{}llrrl@{}}
\toprule
\textbf{Field} & \textbf{Unit} & \textbf{Accepted range} &
\textbf{Reference $r_f$} & \textbf{Type} \\
\midrule
age & years & 18--85 & 55.0 & numeric \\
income\_poverty\_ratio & ratio & 0.0--5.0 & 2.0 & numeric \\
waist & cm & 40--180 & 98.0 & numeric \\
mean\_sbp & mmHg & 70--240 & 125.0 & numeric \\
mean\_dbp & mmHg & 35--140 & 78.0 & numeric \\
glucose & mg/dL & 40--500 & 100.0 & numeric \\
hba1c & \% & 3.0--18.0 & 5.6 & numeric \\
blood\_urea\_nitrogen & mg/dL & 2--150 & 14.0 & numeric \\
uric\_acid & mg/dL & 1.0--15.0 & 5.5 & numeric \\
serum\_albumin & g/dL & 1.5--6.0 & 4.2 & numeric \\
hemoglobin & g/dL & 5.0--20.0 & 14.0 & numeric \\
diabetes\_yes & 0/1 & \{0, 1\} & 0 & binary \\
hypertension\_doctor\_told & 0/1 & \{0, 1\} & 0 & binary \\
taking\_bp\_medicine & 0/1 & \{0, 1\} & 0 & binary \\
\bottomrule
\end{tabular}
\end{table}

The API rejects unknown fields (\texttt{extra="forbid"}) and out-of-range values
with HTTP 422 before the model is invoked. Fields left blank are imputed by the
fitted preprocessing pipeline and counted in $n_{\mathrm{miss}}$.

\chapter{Safety Rules and Output Contract}
\label{app:rules}

\section*{B.1 The nine safety rules}

The following rules appear verbatim in the full system prompt and are mirrored
by the deployed validator.

\begin{enumerate}
  \item Do not state or imply a diagnosis. Never say the person has CKD, kidney
        disease, or kidney failure. This is a screening-based risk level, not a
        diagnosis.
  \item Do not prescribe, recommend, start, stop, change, or dose any medication
        or supplement.
  \item Do not give individualised nutrition restrictions or numeric targets for
        protein, potassium, phosphorus, or sodium. If diet is relevant, suggest
        discussing a personalised plan with a clinician or registered/renal
        dietitian.
  \item Do not replace clinician judgement, \eGFR{} testing, or urine albumin
        (\UACR) testing.
  \item Give wellness education tailored to the patient's specific conditions and
        risk factors shown in the patient profile. Reference their actual
        conditions when giving food, activity, or monitoring advice; do not
        diagnose, prescribe, or treat.
  \item For a moderate or high screening result, always include a reminder to
        confirm the result with a clinician and appropriate laboratory tests.
  \item If the inputs suggest an urgent problem (for example chest pain, severe
        shortness of breath, severe swelling, or very high blood pressure),
        advise seeking prompt medical care.
  \item Base the education on the provided guideline excerpts. If the excerpts do
        not cover something, keep the advice general, safe, and non-specific.
  \item Use calm, supportive, non-alarming, plain language at about an
        eighth-grade reading level.
\end{enumerate}

\section*{B.2 Output contract}

Every response must contain exactly five labelled sections in the following
order, followed by a machine-readable grounding line listing the identifiers of
the guideline excerpts actually used:

\begin{enumerate}
  \item \textbf{Summary} --- names the screening risk level in everyday terms
        and the one or two factors that most influenced it.
  \item \textbf{Food habits} --- eating habits tied to this individual's
        factors; no numeric nutrient targets.
  \item \textbf{Rest and sleep} --- practical sleep and rest habits.
  \item \textbf{Daily routine} --- movement, hydration and monitoring habits
        chosen for this individual's factors.
  \item \textbf{Follow-up reminder} --- when and why to consult a clinician and
        which routine tests to ask about; mandatory for moderate and high.
\end{enumerate}

\chapter{Validator Detection Layers}
\label{app:validator}

\begin{table}[H]
\centering
\caption{Validator layers, severity and disposition.}
\label{tab:app_validator}
\small
\begin{tabularx}{\textwidth}{@{}L{1.35in} X l L{1.05in}@{}}
\toprule
\textbf{Layer} & \textbf{Detects} & \textbf{Severity} & \textbf{Action} \\
\midrule
Structural check & A required section is absent &
--- & Regenerate once, then fall back \\
Diagnosis blocker & Assertion that the person has CKD, kidney disease or kidney
failure & Severe & Safe template \\
Medication blocker & Drug names or classes, dosing expressions, or
start/stop/change instructions & Severe & Safe template \\
Supplement blocker & Recommendation to take a supplement & Severe &
Safe template \\
Nutrition restriction & Numeric individualised protein, potassium, phosphorus or
sodium target & Moderate & Rewrite to dietitian referral \\
Lab-confirmation reminder & Moderate or high result lacking a clinician and
laboratory follow-up mention & Moderate & Append reminder \\
Emergency redirect & Systolic $\geq$180 or diastolic $\geq$120 mmHg, or red-flag
symptom vocabulary & Moderate & Prepend urgent-care advice \\
\bottomrule
\end{tabularx}
\end{table}

Detection is sentence-scoped and negation-aware throughout: a candidate match is
discarded when governed by a negation or an explicit deferral to a clinician,
so that safe guidance containing clinical vocabulary (``do not stop or change
any medicine without talking to your doctor'') is not misclassified as a
violation.

\chapter{Worked Example}
\label{app:example}

The following illustrates a complete request for a high-risk profile. Values are
representative of the benchmark's emergency subgroup.

\section*{D.1 Input}

\begin{quote}\footnotesize\ttfamily\raggedright
age 65,\; waist 112 cm,\; mean\_sbp 186 mmHg,\; mean\_dbp 104 mmHg,\;
glucose 168 mg/dL,\; hba1c 8.5\%,\; blood\_urea\_nitrogen 26 mg/dL,\;
uric\_acid 8.1 mg/dL,\; serum\_albumin 3.6 g/dL,\; hemoglobin 11.9 g/dL,\;
income\_poverty\_ratio 1.2,\; diabetes\_yes 1,\;
hypertension\_doctor\_told 1,\; taking\_bp\_medicine 1
\end{quote}

\section*{D.2 Screening output}

\begin{quote}\footnotesize\ttfamily\raggedright
stage1: p\_atrisk = 0.7447 ($\geq$ 0.44) $\rightarrow$ sent to stage 2\\
stage2: p\_high\_given\_atrisk = 0.7290 ($\geq$ 0.50)\\
risk\_class = high\\
probabilities: low 0.2553, moderate 0.2018, high 0.5428\\
confidence = 0.5428, reliability = moderate\\
top factors: mean\_sbp (raises), hba1c (raises),\\
\hspace*{2em}blood\_urea\_nitrogen (raises), waist (raises)
\end{quote}

\section*{D.3 Retrieval}

Priority topics derived from the risk-raising factors are
\texttt{bp}, \texttt{diabetes}, \texttt{weight}, \texttt{followup}; four
passages are selected, one per topic, drawn from KDIGO 2021 (blood pressure),
the ADA/KDIGO consensus (diabetes), NKF patient education (weight and activity)
and CDC prevention material (follow-up).

\section*{D.4 Validator action}

Because systolic 186 mmHg exceeds the emergency threshold of 180 mmHg, the
urgent-vitals alert block is inserted into the prompt, and the emergency-redirect
layer verifies that urgent-care advice is present in the delivered text,
prepending it otherwise. The disposition for this profile is
\texttt{sanitized}, and the delivered Summary opens with an urgent-care
instruction before restating the screening category.

\chapter{Reproduction Instructions}
\label{app:repro}

\section*{E.1 Environment}

\begin{quote}\footnotesize\ttfamily\raggedright
python -m venv .venv \\
.venv\textbackslash Scripts\textbackslash activate \\
pip install -r requirements.txt \\
pip install -r requirements-llm.txt
\end{quote}

\section*{E.2 Configuration}

\begin{quote}\footnotesize\ttfamily\raggedright
CKD\_LLM\_MODEL=qwen \\
CKD\_LLM\_ADAPTER=llm/finetune/adapters/v1/adapter \\
DATABASE\_URL=postgresql+psycopg2:// \\
\hspace*{2em}user:password@localhost:5432/ckd\_research \\
CKD\_DB\_AUTO\_CREATE=false
\end{quote}

\section*{E.3 Pipeline commands}

\begin{quote}\footnotesize\ttfamily\raggedright
\# database schema \\
python -m alembic upgrade head \\[4pt]
\# build the retrieval index \\
ollama pull bge-m3 \\
python -m llm.rag.embed\_index --embedder ollama-bge-m3 \\[4pt]
\# generate fine-tuning data and train the adapter \\
python -m llm.data.generate\_dataset --n 600 --out-dir phase7\_data \\
python -m llm.finetune.qlora\_train --data-dir phase7\_data \\[4pt]
\# build the evaluation benchmark and run the ablation \\
python -m llm.eval.q1.benchmark \textbackslash \\
\hspace*{2em}--out q1\_benchmark/q1\_profiles\_300.jsonl \\
python -m llm.eval.q1.run\_arms --arm A ... --arm F \\
python -m llm.eval.q1.report --runs q1\_runs --out-dir q1\_report \\[4pt]
\# regenerate thesis figures \\
python thesis/make\_figures.py \\[4pt]
\# run the application \\
uvicorn backend.main:app --reload \\
cd frontend \&\& npm run dev
\end{quote}

\section*{E.4 Verification}

The health endpoint confirms which model is actually serving requests:

\begin{quote}\footnotesize\ttfamily\raggedright
GET /generate-guidance/health \\[3pt]
\{ \\
\hspace*{1.5em}"model\_spec": \\
\hspace*{3em}"qwen:Qwen/Qwen3-4B-Instruct-2507+lora:v1/adapter", \\
\hspace*{1.5em}"adapter\_active": true, \\
\hspace*{1.5em}"prompt": "compact", \\
\hspace*{1.5em}"mode": "production" \\
\}
\end{quote}
